\documentclass[10pt,a4paper]{amsart}
\usepackage{amsmath,amssymb,amsthm,graphicx}
% Retained minimal preamble for the Stacks Project chapter more-algebra.tex (More on Algebra).
% The chapter file pinned-original-source.tex carries only \input{preamble} on its first line, so no
% preamble text lives inside the chapter itself.  The Stacks Project uses one single shared preamble for
% every chapter; the snapshot stages it identically as sources/stacks-chow/preamble.tex and
% sources/stacks-etale-cohomology/preamble.tex.  The block below is copied mechanically, line for line,
% from sources/stacks-chow/preamble.tex lines 14-18 (the Xy-pic setup) and 154-217 (the native theorem
% declarations together with the mathematical macros).  Nothing was added, removed or redefined; the
% wrapper class and the amsmath/amssymb/amsthm/graphicx base packages are supplied by the assembler.
% Chapter: more-algebra.tex -- More on Algebra.
% For commutative diagrams we use Xy-pic
\usepackage[all]{xy}

% We use 2cell for 2-commutative diagrams.
\xyoption{2cell}

% Theorem environments.
%
\theoremstyle{plain}
\newtheorem{theorem}[subsection]{Theorem}
\newtheorem{proposition}[subsection]{Proposition}
\newtheorem{lemma}[subsection]{Lemma}

\theoremstyle{definition}
\newtheorem{definition}[subsection]{Definition}
\newtheorem{example}[subsection]{Example}
\newtheorem{exercise}[subsection]{Exercise}
\newtheorem{situation}[subsection]{Situation}

\theoremstyle{remark}
\newtheorem{remark}[subsection]{Remark}
\newtheorem{remarks}[subsection]{Remarks}

\numberwithin{equation}{subsection}

% Macros
%
\def\lim{\mathop{\mathrm{lim}}\nolimits}
\def\colim{\mathop{\mathrm{colim}}\nolimits}
\def\Spec{\mathop{\mathrm{Spec}}}
\def\Hom{\mathop{\mathrm{Hom}}\nolimits}
\def\Ext{\mathop{\mathrm{Ext}}\nolimits}
\def\SheafHom{\mathop{\mathcal{H}\!\mathit{om}}\nolimits}
\def\SheafExt{\mathop{\mathcal{E}\!\mathit{xt}}\nolimits}
\def\Sch{\mathit{Sch}}
\def\Mor{\mathop{\mathrm{Mor}}\nolimits}
\def\Ob{\mathop{\mathrm{Ob}}\nolimits}
\def\Sh{\mathop{\mathit{Sh}}\nolimits}
\def\NL{\mathop{N\!L}\nolimits}
\def\CH{\mathop{\mathrm{CH}}\nolimits}
\def\proetale{{pro\text{-}\acute{e}tale}}
\def\etale{{\acute{e}tale}}
\def\QCoh{\mathit{QCoh}}
\def\Ker{\mathop{\mathrm{Ker}}}
\def\Im{\mathop{\mathrm{Im}}}
\def\Coker{\mathop{\mathrm{Coker}}}
\def\Coim{\mathop{\mathrm{Coim}}}

% Boxtimes
%
\DeclareMathSymbol{\boxtimes}{\mathbin}{AMSa}{"02}

%
% Macros for moduli stacks/spaces
%
\def\QCohstack{\mathcal{QC}\!\mathit{oh}}
\def\Cohstack{\mathcal{C}\!\mathit{oh}}
\def\Spacesstack{\mathcal{S}\!\mathit{paces}}
\def\Quotfunctor{\mathrm{Quot}}
\def\Hilbfunctor{\mathrm{Hilb}}
\def\Curvesstack{\mathcal{C}\!\mathit{urves}}
\def\Polarizedstack{\mathcal{P}\!\mathit{olarized}}
\def\Complexesstack{\mathcal{C}\!\mathit{omplexes}}
% \Pic is the operator that assigns to X its picard group, usage \Pic(X)
% \Picardstack_{X/B} denotes the Picard stack of X over B
% \Picardfunctor_{X/B} denotes the Picard functor of X over B
\def\Pic{\mathop{\mathrm{Pic}}\nolimits}
\def\Picardstack{\mathcal{P}\!\mathit{ic}}
\def\Picardfunctor{\mathrm{Pic}}
\def\Deformationcategory{\mathcal{D}\!\mathit{ef}}
\usepackage{cleveref}
\crefname{originallemma}{original lemma}{original lemmas}
\Crefname{originallemma}{Original lemma}{Original lemmas}
\makeatletter
\newlabel{more-algebra-lemma-descent-fs}{{07EH}{source}{}{sourcepointer-1}{}}
\newlabel{more-algebra-lemma-descent-fs@cref}{{[originallemma][0][]07EH}{[1][0][]source}{}{}{}}
\makeatother
\crefname{originallemma}{original lemma}{original lemmas}
\Crefname{originallemma}{Original lemma}{Original lemmas}
\makeatletter
\newlabel{definition-formally-smooth}{{07EB}{source}{}{sourcepointer-2}{}}
\newlabel{definition-formally-smooth@cref}{{[originallemma][0][]07EB}{[1][0][]source}{}{}{}}
\makeatother
\title{Formal smoothness descends along a ring map that is a direct summand}
\author{The Stacks Project Authors}
\date{stacks-project:more-algebra.tex:local-snapshot-20261012; GNU FDL 1.2 or later}
\begin{document}
\maketitle
\setcounter{subsection}{1}
\noindent\textit{Editorial note.} All delivered mathematical text is the author's own, copied from the pinned Stacks Project chapter file more-algebra.tex; no statement or proof step was rewritten, reordered or removed.\par\medskip
\noindent\textit{Editorial note.} The chapter file pinned-original-source.tex carries only \textbackslash{}input\{preamble\} on its first line, so no preamble text lives inside the chapter. The retained file sources/stacks-more-algebra/retained-preamble.tex is a mechanical, line-for-line copy of the project's single shared preamble as staged in the snapshot (Xy-pic setup plus theorem declarations and mathematical macros). Nothing was added, removed or redefined, and the excerpt serves only as preamble, never as a span coordinate.\par\medskip
\noindent\textit{Editorial note.} Both counted spans start at the \textbackslash{}begin line of their environment, so the statement and the author's proof are delivered whole; surrounding prose of the section is deliberately outside the counted spans and is carried as standard prerequisites instead.\par\medskip
\noindent\textit{Editorial note.} The citation fields (source\_id, source\_version, authors, source\_url, license.name, source\_title) are identical to those of every other card this lane has produced from this chapter, so the batch selector treats them as one source.\par\medskip
\noindent\textit{Allowed standard prerequisites (editorial).}\begin{itemize}
\item Formal smoothness in the adic topology, defined through the lifting of diagrams along square-zero ideals. the author cites Definition definition-formally-smooth at proof lines 9389-9392. Supplies the property to be descended and the form of the diagram that has to be lifted.
\item Base change of a lifting diagram along the ring map R to R'. used at proof lines 9393-9400. Produces the diagram to which the hypothesis applies.
\item A direct summand decomposition R' = R plus C and the resulting decomposition of the base changed target, with the complement part an ideal because the relevant ideal squares to zero. used at proof lines 9400-9416. Is the computation that produces the arrow R to S to A required by the definition.
\end{itemize}\medskip
\ifdefined\hypertarget\hypertarget{sourcepointer-1}{}\fi
\noindent\textit{Source reference.} Original lemma 07EH: focused-seeds/AGP/sources/stacks-snapshot-20261011/tags.txt, source lines 9447--9447. The lemma's own canonical tag, quoted so that readers can resolve the label inside the delivered text; it points at the counted statement itself and carries no proof content.\par\medskip
\ifdefined\hypertarget\hypertarget{sourcepointer-2}{}\fi
\noindent\textit{Source reference.} Original lemma 07EB: focused-seeds/AGP/sources/stacks-snapshot-20261011/tags.txt, source lines 9441--9441. Navigation pointer to a statement the counted proof refers to but does not contain; nothing load-bearing is delegated to it.\par\medskip
\begin{lemma}
\label{lemma-descent-fs}
Let $R$, $S$ be rings. Let $\mathfrak n \subset S$ be an ideal.
Let $R \to R'$ be a ring map. Set $S' = S \otimes_R R'$ and
$\mathfrak n' = \mathfrak nS$. If
\begin{enumerate}
\item the map $R \to R'$ embeds $R$ as a direct summand of $R'$
as an $R$-module, and
\item $R' \to S'$ is formally smooth for the $\mathfrak n'$-adic topology,
\end{enumerate}
then $R \to S$ is formally smooth in the $\mathfrak n$-adic topology.
\end{lemma}

\begin{proof}
Let a solid diagram
$$
\xymatrix{
S \ar[r] & A/J \\
R \ar[u] \ar[r] & A \ar[u]
}
$$
as in Definition \ref{definition-formally-smooth} be given.
Set $A' = A \otimes_R R'$ and $J' = \Im(J \otimes_R R' \to A')$.
The base change of the diagram above is the diagram
$$
\xymatrix{
S' \ar[r] \ar@{-->}[rd]^{\psi'} & A'/J' \\
R' \ar[u] \ar[r] & A' \ar[u]
}
$$
with continuous arrows. By condition (2) we obtain the dotted arrow
$\psi' : S' \to A'$. Using condition (1) choose a direct summand decomposition
$R' = R \oplus C$ as $R$-modules. (Warning: $C$ isn't an ideal in $R'$.)
Then $A' = A \oplus A \otimes_R C$. Set
$$
J'' = \Im(J \otimes_R C \to A \otimes_R C) \subset J' \subset A'.
$$
Then $J' = J \oplus J''$ as $A$-modules. The image of the composition
$\psi : S \to A'$ of $\psi'$ with $S \to S'$ is contained in
$A + J' = A \oplus J''$. However, in the ring $A + J' = A \oplus J''$
the $A$-submodule $J''$ is an ideal! (Use that $J^2 = 0$.) Hence the
composition $S \to A + J' \to (A + J')/J'' = A$ is the arrow we were
looking for.
\end{proof}

\end{document}
